% Figure: mechanism occurrences by claim status (Table A3). Thesis palette: darker indigo = stronger evidence.
\begin{tikzpicture}
\begin{axis}[ntaxis, xbar stacked, width=9.2cm, height=6.6cm, bar width=8pt,
  xmin=0, xmax=26, xtick={0,5,...,25}, xlabel={Occurrences (of 130)},
  ytick={0,...,11}, yticklabel style={align=right,text width=6.2cm},
  yticklabels={{Indicator meaningless as defined},{Statistical misspecification},
    {Comparability between installations},{Other},{Unmeasured quantity},{Base or non-productive load},
    {Missing, wrong or discontinuous data},{Non-stationarity of the relationship},
    {Allocation between products},{Baseline drift or obsolescence},{Validity domain exceeded},{Series too short}},
  legend style={font=\scriptsize, draw=none, at={(1,0.02)}, anchor=south east, cells={anchor=west}},
  legend cell align=left,
  after end axis/.append code={
% segment labels (centre of each segment) and totals
\node[font=\scriptsize,text=white] at (axis cs:3,0) {6};
\node[font=\scriptsize,text=corTinta] at (axis cs:10,0) {8};
\node[font=\scriptsize,text=corTinta] at (axis cs:19,0) {10};
\node[font=\scriptsize,text=corTinta,anchor=west] at (axis cs:24,0) {24};
\node[font=\scriptsize,text=white] at (axis cs:2.5,1) {5};
\node[font=\scriptsize,text=corTinta] at (axis cs:6.5,1) {3};
\node[font=\scriptsize,text=corTinta] at (axis cs:14.5,1) {13};
\node[font=\scriptsize,text=corTinta,anchor=west] at (axis cs:21,1) {21};
\node[font=\scriptsize,text=white] at (axis cs:1,2) {2};
\node[font=\scriptsize,text=corTinta] at (axis cs:5.5,2) {7};
\node[font=\scriptsize,text=corTinta] at (axis cs:12,2) {6};
\node[font=\scriptsize,text=corTinta,anchor=west] at (axis cs:15,2) {15};
\node[font=\scriptsize,text=white] at (axis cs:1.5,3) {3};
\node[font=\scriptsize,text=corTinta] at (axis cs:7.5,3) {9};
\node[font=\scriptsize,text=corTinta] at (axis cs:12.5,3) {1};
\node[font=\scriptsize,text=corTinta,anchor=west] at (axis cs:13,3) {13};
\node[font=\scriptsize,text=white] at (axis cs:0.5,4) {1};
\node[font=\scriptsize,text=corTinta] at (axis cs:5.5,4) {9};
\node[font=\scriptsize,text=corTinta,anchor=west] at (axis cs:10,4) {10};
\node[font=\scriptsize,text=white] at (axis cs:2.5,5) {5};
\node[font=\scriptsize,text=corTinta] at (axis cs:7.5,5) {5};
\node[font=\scriptsize,text=corTinta,anchor=west] at (axis cs:10,5) {10};
\node[font=\scriptsize,text=white] at (axis cs:1.5,6) {3};
\node[font=\scriptsize,text=corTinta] at (axis cs:5.5,6) {5};
\node[font=\scriptsize,text=corTinta] at (axis cs:8.5,6) {1};
\node[font=\scriptsize,text=corTinta,anchor=west] at (axis cs:9,6) {9};
\node[font=\scriptsize,text=white] at (axis cs:0.5,7) {1};
\node[font=\scriptsize,text=corTinta] at (axis cs:3.5,7) {5};
\node[font=\scriptsize,text=corTinta] at (axis cs:7,7) {2};
\node[font=\scriptsize,text=corTinta,anchor=west] at (axis cs:8,7) {8};
\node[font=\scriptsize,text=corTinta] at (axis cs:1.5,8) {3};
\node[font=\scriptsize,text=corTinta] at (axis cs:5.5,8) {5};
\node[font=\scriptsize,text=corTinta,anchor=west] at (axis cs:8,8) {8};
\node[font=\scriptsize,text=white] at (axis cs:0.5,9) {1};
\node[font=\scriptsize,text=corTinta] at (axis cs:1.5,9) {1};
\node[font=\scriptsize,text=corTinta] at (axis cs:4.5,9) {5};
\node[font=\scriptsize,text=corTinta,anchor=west] at (axis cs:7,9) {7};
\node[font=\scriptsize,text=white] at (axis cs:0.5,10) {1};
\node[font=\scriptsize,text=corTinta] at (axis cs:2,10) {2};
\node[font=\scriptsize,text=corTinta] at (axis cs:3.5,10) {1};
\node[font=\scriptsize,text=corTinta,anchor=west] at (axis cs:4,10) {4};
\node[font=\scriptsize,text=corTinta] at (axis cs:0.5,11) {1};
\node[font=\scriptsize,text=corTinta,anchor=west] at (axis cs:1,11) {1};
}]
\addplot[draw=corDestaque, fill=corDestaque] coordinates
  {(6,0) (5,1) (2,2) (3,3) (1,4) (5,5) (3,6) (1,7) (0,8) (1,9) (1,10) (0,11)};
\addplot[draw=corDestaque, fill=corDestaqueC] coordinates
  {(8,0) (3,1) (7,2) (9,3) (9,4) (0,5) (5,6) (5,7) (3,8) (1,9) (2,10) (0,11)};
\addplot[draw=corCinza, fill=corCinzaC] coordinates
  {(10,0) (13,1) (6,2) (1,3) (0,4) (5,5) (1,6) (2,7) (5,8) (5,9) (1,10) (1,11)};
\legend{Demonstrated, Applied, General}
\end{axis}
\end{tikzpicture}
